% Page budget: 1.35 pages | Owns: gantry coordinates, first-principles equations, scheduling, LPV-LFR realisation, and physical parameterisation.
% WRITING GUIDE
% Purpose: Define the physical coordinates and derive the exact continuous-time LPV-LFR baseline used by every later section.
% Primary reference for Section II-B and its appendix material:
% docs/Thesis-documentation/LPV_LFR_Stepwise_Derivation/main.tex
% (the version shared with Jasper, including the G-Delta interconnection figure).
% Follow its signal definitions, direct mass-equation construction, block partitioning,
% loop-closing verification, and allocation of detailed algebra to the appendix.
% Supporting checks/details only: LPV_LFR_internal_loop_Well_posedness,
% LPV_Mass_invertible, LPV_LFR_Justification_Drenth, LPV_LFR_Derivation,
% and docs/fp-model-structure.md. Where their presentation differs, follow
% LPV_LFR_Stepwise_Derivation for Section II-B.
% Sources: README "Source map per subsection", rows II-A, II-B, II-C and II Frozen-LTI baseline; mandatory papers and the reference gate as in the README.
% Research-plan anchor: Preserve Aspect 1's explicit position dependence; update the original discretisation wording to the final method.
% Current decisions: Continuous-time physics, endogenous scheduler Y, fixed-step RK4, and identifiable parameter combinations.
% Choices to justify: Coordinate frame, Y as scheduler, continuous-time formulation, RK4 step and substeps, full six-channel LFR realization, and reduced physical parameterization. Give the physical or numerical reason for each choice, not only its definition.
% Cautions: The final pipeline uses the full six-channel scheduling loop, not the experimental six-to-four SVD reduction. Resolve the supervisor comment about dual-gantry terminology; do not copy unverified bibliography entries.
% Planned figures: Physical coordinate schematic and compact interconnection between G and Delta(Y).
% Section is finished when: The LFR collapses to the original model, well-posedness assumptions are stated, and notation matches Section 03.
%
% LPV-LFR DERIVATION ROADMAP
% 1. Define q=[X,Theta,Y]^T, the force transformation, measured outputs, assumptions, and the admissible range of Y.
% 2. State M(Y)qdd+Cqd+Kq=f and expose the structural dependence M(Y)=M0+Y M1+Y^2 M2.
% 3. Construct the LFR directly from the polynomial mass equation:
%       a1=Y qdd, a2=Y a1, M0 qdd + M1 a1 + M2 a2 = f_net,
%       z=[qdd; a1], w=[a1; a2] (stepwise latent names; qdd is not renamed to a).
% 4. Collect w=Delta(Y)z with Delta(Y)=Y I6 and give the constant block-level G realization with all signal dimensions.
% 5. Eliminate w and z in a short exactness check to recover the original equation of motion.
% 6. State the well-posedness condition over the complete scheduling range and relate it to invertibility of M(Y); place the longer proof outside the main text.
% 7. Only after the conceptual construction, mention that the implementation
%    evaluates the resulting rational expression in Horner form. Put
%    M(Y)^(-1)=[N0+Y N1+Y^2 N2]/d(Y), the explicit Ni entries, and the
%    determinant coefficients in the appendix or technical supplement.
% 8. Keep this derivation continuous-time; Section III-A owns the RK4 scheme, Section V only its step size.
%
% DERIVATION WARNING
% Do not start from v=M(Y)^(-1)f_net. Follow LPV_LFR_Stepwise_Derivation:
% construct the repeated scheduling loop from qdd, a1=Y qdd, a2=Y a1 with
% z=[qdd; a1], w=[a1; a2], then verify it by closing the loop. Older inverse-first notes
% may be used only as supporting algebra. Verify terminology against the
% primary literature by Drenth and Hoekstra.
% MUST ESTABLISH (II-B after eq:delta, agreed 2026-10-06):
% 1. G is a constant LTI system in Drenth Eq. (6) form; its first block row is the state equation.
% 2. Closing the algebraic loop gives back eq:eom exactly.
% 3. Well posed iff M(Y) nonsingular, true for every Y; hence closed-form M(Y)^{-1}=N(Y)/d(Y)
%    (N_i, d(Y) and Horner evaluation in the appendix). The reduced-parameter bound stays in II-C.
% MUST ESTABLISH (II-C, agreed 2026-10-07):
% 1. The fourteen raw parameters enter M0,M1,M2,C,K only through theta_base (eq:phi); this work's own.
% 2. Argument in the main text, ordered toward the conclusion: input-output behaviour fixes the matrices;
%    twelve entries, three equal m_h, ten independent; so exactly ten combinations are structurally
%    identifiable; only then named theta_base (eq:phi). Model-only (no data), no new symbols, no
%    Proposition/proof environment (agreed 2026-10-07).
% 3. Four lost directions in one sentence; estimates reported in theta_base.
% 4. Positive theta_base does not give M(Y)>0; that holds iff eq:lfr_admissibility (proof in app:lfr-wellposed).
% 5. One sentence: training keeps eq:lfr_admissibility at every iterate. The m_Delta map (eq:mdelta_map)
%    and its justification are stated in Section III-C (agreed 2026-10-07), not here.
% Compact form agreed 2026-10-07: no separate definition paragraph, no null-direction display.
% Moved to Section V as bullets: rho value, combination-error metric with the m_Delta normalisation, raw vs
% reduced coordinates per arm, practical identifiability on the records. Appendix app:params keeps the table only.
% Main text: assumptions, polynomial dependence, direct LFR construction, compact G, exactness, well-posedness statement, and one implementation sentence.
% Appendix/supplement: expanded M_i and N_i matrices, determinant polynomial, longer invertibility/well-posedness proof, and symbolic verification.
% MUST ESTABLISH (II-A, PROPOSED cycle 09, inferred from README ownership and source-map row II-A):
% 1. The system is the dual-drive gantry of Garcia-Herreros et al.; q=(X,Theta,Y) avoids a closed kinematic
%    chain and exposes the load-distribution coupling (Garcia Sec. 2.2); P links q to the measured frame.
% 2. eq:eom under Garcia's small-angle and no-Coriolis simplifications (Sec. 2.3).
% 3. Adapted: all three coordinates kept (Garcia's (11) keeps X, Theta); Coulomb terms of Eq. (6) omitted with
%    a todo naming D-022 vs D-204 (source-map row II-A); L_b fixed because it defines q.
% MUST ESTABLISH (II-D compared baselines, PROPOSED cycle 09, from README ownership, source-map row
%    "II Frozen-LTI baseline", THESIS-RESULTS step 1): M_LPV and M_LTI named together with ports u_act to
%    q_act; both untrained at the true parameters; M_LTI evaluates Delta at a fixed Y_op (value in Section V).
% CYCLE 09 NOTES (header only):
% - Corrections to the blocks (skill 2a, claims kept): II-B item 1 cites Drenth Eq. (6), a discrete-time form
%   whose LTI part is named M; the CT form is cited from drenth2025thesis Eq. (2.1) (source-map row II-B).
%   G is displayed in full in the main text with its ports u_act and q_act (README Math standard item 1;
%   Derivation policy item 5: the well-posedness claim relies on D_zw). II-C item 4: theta_base contains the
%   signed m_Delta, so "positive theta_base" is written as "m_h, m_Sigma, J_eff > 0" (code: m_diff signed).
%   II-B item 3 "N(Y)/d(Y)": d is the payload offset, so the appendix writes adj M(Y)/det M(Y).
% - Reasons and sources: coordinates (Garcia Sec. 2.2, citation-log); small angle and Coriolis (Garcia
%   Sec. 2.3); Y as scheduler and known scheduling map (D-004; Drenth 2025 Sec. 3.2 contrast); continuous time
%   (D-018, D-020: keep the physical parameter structure, RK4 at runtime, no pre-discretisation); LFR as the
%   baseline form of LPV-LFR augmentation (drenth2025thesis Sec. 5.2, Eq. (5.1), read 2026-10-09);
%   implementation equivalence (FS/blocks.py::Gantry_State_Block._deriv_with: closed-form N(Y)/d(Y) in Horner
%   form, z, w, then xdot through A_combined = [Ax Bw Bu] with P applied to the stage forces).
% - Own math checked symbolically this session (sympy): M split; det(I6 - D_zw Delta(Y)) = det M(Y)/det M0;
%   twelve nonzero entries, three equal to +-m_h, rank ten; completed square and raw-parameter identity of
%   app:lfr-wellposed.
% - Left to other sections: RK4 step and substeps (III-A scheme, V value); Y_op value, parameters of the
%   truth, stroke and data ranges (V); m_Delta map (III-C); SVD reduction excluded (header Caution, meeting
%   audit candidate-thesis-impact.md).
% - Resolved original todos: "Does this make sense to put here" (owned by Section III opening); coordinate
%   reason and "nonuniform ... moving payload" (Garcia Sec. 2.2 in prose); "geometric dimensions" (L_b, d in
%   the where clause); "Look at this section" (rewrite); Telica remark (simulation only, THESIS-RESULTS Sec. 2);
%   operating range of Y (M(Y)>0 for every Y in II-B; ranges in V); "source for self-scheduled" (Toth
%   Def. 7.2, Drenth Sec. 3.2); SVD reduction (excluded); "Left here" (draft complete). The Coulomb insertion
%   todo is replaced by the D-022 vs D-204 todo; its Hoekstra Cond. 8 candidate is unverified and dropped.
% - Removed displays eq:lfr_mass_factorization, eq:lfr_chain (content in eq:Msplit, eq:lfr_constant). New
%   labels sec:system, sec:baselines, eq:wellposed, eq:baselines. 99_appendix.tex app:lfr-details updated
%   (partition now in the main text; d(Y) renamed det M(Y); two cut remarks now in the main text), and
%   app:lfr-wellposed (a renamed beta, a clashed with a_1).
% - Stale in other files (not edited): 01_introduction.tex Aspect 1 todo ("LPV versus frozen LTI exists only as
%   a baseline check") is now covered by eq:baselines; Introduction writes "dual-gantry" (terminology todo).
\section{Dual-Drive Gantry System and Physics Baseline}
\ifdraft
\begin{itemize}
\item Opening: derives the baseline that Section~III augments, the first contribution of Section~I; names II-A to II-D (style profile)
\end{itemize}
\fi

This section derives the baseline that Section~\ref{sec:augmentation}
augments: a continuous-time LPV-LFR realisation of the dual-drive gantry
with rational dependence on the payload position, the first contribution of
Section~I.
Section~\ref{sec:system} states the equations of motion,
Section~\ref{sec:lfr} their LPV-LFR realisation and its well-posedness,
Section~\ref{sec:params} the parameter combinations that input-output data
identify, and Section~\ref{sec:baselines} the two baselines that the results
compare.
\todo{Missing: contribution wording. Section~\ref{sec:params} also delivers the ten identifiable combinations and their admissibility condition, which the first contribution does not name, and ``compact'' sits beside a realisation that is not claimed minimal. Candidate: ``Derive an exact continuous-time LPV-LFR realisation of the dual-drive gantry baseline with rational dependence on payload position, parameterised by its ten identifiable combinations'' (D-191, D-229).}

\subsection{System and first-principles baseline}\label{sec:system}
\ifdraft
\begin{itemize}
\item Two X drives carry a cross-arm through flexible joints; a Y drive moves the payload along it (Garcia Sec.~2.2)
\item $q=(X,\Theta,Y)$ avoids a closed kinematic chain and exposes the load-distribution coupling; $P$ links it to the measured frame (Garcia Sec.~2.2; adopted)
\item Equations of motion with $\cos\Theta\approx1$, $\sin\Theta\approx0$ and no Coriolis term (Garcia Sec.~2.3; adopted)
\item Adapted: all three coordinates kept, unlike Garcia's (11); Coulomb terms of Garcia Eq.~(6) omitted (D-022 vs D-204 (?)); $L_b$ fixed because it defines $q$
\item Nominal values from Garcia Table~I, in the appendix
\end{itemize}
\fi
The system is the dual-drive gantry of Garc\'ia-Herreros et
al.~\cite{garcia2013model} (Fig.~\ref{fig:gantry_system}).
Two parallel linear drives carry a cross-arm through flexible joints.
A third drive moves the payload along the cross-arm.
\todo{Missing: the supervisor comment on dual-gantry terminology (header Caution) has no source in the decisions or the meeting audit. Candidate: ``dual-drive gantry'', the term of the title of~\cite{garcia2013model} and of the research plan; the Introduction still writes ``dual-gantry''.}
\begin{figure}[b]
  \centering
  \def\gantryabsorber{0}% baseline only; the absorber version is in Section III
  \input{figures/gantry_system_with_absorber}
  \caption{Lumped-parameter gantry schematic adapted from the mechanical
  arrangement in \cite{garcia2013model}. Red labels denote actuator coordinates
  and blue labels generalised coordinates.}
  \label{fig:gantry_system}
\end{figure}

We model the gantry in the generalised coordinates of Garc\'ia-Herreros et
al., which avoid a closed kinematic chain and expose the coupling caused by
the non-uniform load distribution over the two
drives~\cite[Sec.~2.2]{garcia2013model}.
The drives act and are measured in the actuator frame, related to these
coordinates by
\begin{equation}
  q = P^{-\top}\qact,
  \qquad
  u = P\uact,
  \qquad
  P =
  \begin{bmatrix}
    1 & 1 & 0\\
    \tfrac{L_b}{2} & -\tfrac{L_b}{2} & 0\\
    0 & 0 & 1
  \end{bmatrix},
\label{eq:Ptransform}
\end{equation}
where $\qact=[X_1\;X_2\;Y]^\top\in\R^3$ and
$\uact=[F_{X1}\;F_{X2}\;F_Y]^\top\in\R^3$ are the measured positions and the
forces of the drives, $q=[X\;\Theta\;Y]^\top\in\R^3$ holds the cross-arm
translation $X$, its yaw angle $\Theta$ and the payload position $Y$ measured
from the centre of the cross-arm, $u\in\R^3$ is the generalised force, and
$L_b$ is the cross-arm length between the flexible joints.

Following Garc\'ia-Herreros et al., we set $\cos\Theta\approx1$ and
$\sin\Theta\approx0$, since $\Theta$ stays within tens of microradians in
operation~\cite[Sec.~2.3]{garcia2013model}.
As they do, we neglect the Coriolis and centripetal terms.
The equations of motion are then
\begin{equation}
  M(Y)\ddot{q}(t) + C\dot{q}(t) + K q(t) = u(t),
  \label{eq:eom}
\end{equation}
{\footnotesize
\setlength{\arraycolsep}{3pt}
\begin{equation}
M(Y)=
\begin{bmatrix}
m_1+m_2+m_b+m_h & \tfrac{L_b}{2}(m_1-m_2)-m_hY & 0\\
\tfrac{L_b}{2}(m_1-m_2)-m_hY &
\substack{J_b+J_h+\tfrac{L_b^2}{4}(m_1+m_2)\\{}+m_h(d^2+Y^2)} & -m_hd\\
0 & -m_hd & m_h
\end{bmatrix},
\label{eq:mass_matrix}
\end{equation}
\begin{equation}
\begin{aligned}
C&=
\begin{bmatrix}
c_{g1}+c_{g2} & \tfrac{L_b}{2}(c_{g1}-c_{g2}) & 0\\
\tfrac{L_b}{2}(c_{g1}-c_{g2}) &
c_{b1}+c_{b2}+\tfrac{L_b^2}{4}(c_{g1}+c_{g2}) & 0\\
0 & 0 & c_y
\end{bmatrix},\\[1ex]
K&=
\begin{bmatrix}
0&0&0\\
0&k_{b1}+k_{b2}&0\\
0&0&0
\end{bmatrix},
\end{aligned}
\label{eq:damping_stiffness_matrices}
\end{equation}
}
where $M:\R\to\R^{3\times3}$ is the inertia matrix, $C,K\in\R^{3\times3}$
are the damping and stiffness matrices, $m_1$, $m_2$, $m_b$ and $m_h$ are the
masses of the two X drives, the cross-arm and the payload, $J_b$ and $J_h$ the
yaw inertias of the cross-arm and the payload, $d$ the offset of the payload
from the cross-arm, $c_{g1}$, $c_{g2}$ and $c_y$ the damping of the drives,
$c_{b1}$, $c_{b2}$, $k_{b1}$ and $k_{b2}$ the damping and stiffness of the
flexible joints, and $\theta\in\R^{14}$ collects these fourteen parameters,
with nominal values from~\cite[Table~I]{garcia2013model} in
Appendix~\ref{app:params}.

The baseline keeps more coordinates and fewer forces than the control model
of Garc\'ia-Herreros et al.
Their simplified model keeps only $X$ and
$\Theta$~\cite[Eq.~(11)]{garcia2013model}.
We keep all three coordinates, because $Y$ schedules $M(Y)$ and the data
hold all three measured positions.
We omit the Coulomb friction terms of their generalised
forces~\cite[Eq.~(6)]{garcia2013model}.
The length $L_b$ is fixed and not part of $\theta$, because it defines the
coordinates $q$ in \eqref{eq:Ptransform}.
\todo{Missing: reason for omitting the Coulomb terms. D-022 keeps the baseline as Garc\'ia-Herreros et al.\ derived it, which includes them; D-204 keeps friction in the truth only, so that it stays an effect the augmentation must capture (its Karnopp law superseded by the tanh law of D-224, D-226); no entry states why the baseline omits it and Section~\ref{sec:setup} does not either (README source map row II-A). Candidate: state the D-204 reason here.}

\subsection{Self-scheduled LPV-LFR representation}\label{sec:lfr}
\ifdraft
\begin{itemize}
\item $Y$ is a state and the scheduling variable (quasi-LPV, T\'oth Def.~7.2; self-scheduled, Drenth Sec.~3.2); its scheduling map is known, unlike Drenth's jointly estimated map (adapted); $M(Y)^{-1}$ makes the model rational in $Y$, which an LPV-LFR represents exactly (Drenth thesis Sec.~2.1)
\item $M(Y)=M_0+YM_1+Y^2M_2$; $Y$ enters only through $Y\ddot q$ and $Y^2\ddot q$, realised by repeating $Y$, which keeps the coupling of $Y$ and $Y^2$ (this work; precedent Drenth thesis Eq.~(2.9), Sec.~2.1.1)
\item Latent variables $z,w$ and Delta-block $\Dsch(Y)=YI_6$, the full six channels; not claimed minimal (this work; reason (?))
\item LTI part $G$ in the continuous-time form of Drenth thesis Eq.~(2.1), displayed with its ports $\uact$ and $\qact$; constant in $Y$ (adapted)
\item Closing the algebraic loop gives back \eqref{eq:eom} exactly (this work)
\item Well posed iff $M(Y)$ nonsingular, true for every $Y$; the implementation solves the loop in closed form, and the LFR stays the baseline form of LPV-LFR augmentation (Drenth thesis Eq.~(2.4), Sec.~5.2; this work)
\item Continuous time, unlike Drenth et al.; Section~III-A discretises it (D-018, D-020)
\end{itemize}
\fi
The payload position enters \eqref{eq:eom} only through $M(Y)$.
Since $Y$ is a component of the state, the baseline is a quasi-LPV model
with scheduling variable
$Y$~\cite[Def.~7.2]{toth2010modeling}, which Drenth et al.\ call
self-scheduled~\cite[Sec.~3.2]{drenth2025lpvlfr}.
Unlike their self-scheduled models, whose scheduling map is estimated jointly
with the model, ours is known: $Y$ is the
third entry of $q$.
Solving \eqref{eq:eom} for $\ddot q$ requires $M(Y)^{-1}$, whose entries are
rational in $Y$.
An LPV-LFR represents such rational dependence exactly, as a constant LTI
system $G$ in feedback with a block $\Dsch$ that is linear in the scheduling
variable~\cite[Sec.~2.1]{drenth2025thesis} (Fig.~\ref{fig:lfr}).
We construct this realisation for the gantry from the polynomial structure of
$M(Y)$.

\begin{figure}[t]
  \centering
  \input{figures/lfr_loop}
  \caption{Self-scheduled baseline interconnection. The Delta-block closes
  the latent variables of $G$ through $w=\Dsch(Y)z$, with
  $\Dsch(Y)=YI_6$ and $Y$ a component of the state.}
  \label{fig:lfr}
\end{figure}

The inertia matrix \eqref{eq:mass_matrix} is quadratic in $Y$,
\begin{equation}
  M(Y)=M_0+YM_1+Y^2M_2,
  \label{eq:Msplit}
\end{equation}
where $M_0=M(0)$ and $M_1,M_2\in\R^{3\times3}$ depend on $m_h$ only
(Appendix~\ref{app:lfr-details}).
$Y$ therefore enters \eqref{eq:eom} only through $Y\ddot q$ and
$Y^2\ddot q=Y(Y\ddot q)$.

As in the rational embedding of Drenth~\cite[Eq.~(2.9)]{drenth2025thesis},
we realise $Y^2$ by applying $Y$ twice, not as a second scheduling variable,
because independent scheduling variables would admit combinations of $Y$ and
$Y^2$ that cannot occur~\cite[Sec.~2.1.1]{drenth2025thesis}.
With the latent signals $a_1=Y\ddot q$ and $a_2=Ya_1$, \eqref{eq:eom}
becomes
\begin{equation}
  M_0\ddot q+M_1a_1+M_2a_2=u-C\dot q-Kq,
  \label{eq:lfr_constant}
\end{equation}
where $a_1,a_2\in\R^3$.
The scheduling is collected in the Delta-block
\begin{equation}
  w=\begin{bmatrix}a_1\\a_2\end{bmatrix}
   =\underbrace{\begin{bmatrix}YI_3&0\\0&YI_3\end{bmatrix}}_{\Dsch(Y)=YI_6}
    \begin{bmatrix}\ddot q\\a_1\end{bmatrix}
   =\Dsch(Y)z,
  \label{eq:delta}
\end{equation}
where $z,w\in\R^6$ are the latent variables and
$\Dsch:\R\to\R^{6\times6}$.

Solving \eqref{eq:lfr_constant} for $\ddot q$ with the constant $M_0^{-1}$
and substituting \eqref{eq:Ptransform} gives the LTI part, in the
continuous-time form of~\cite[Eq.~(2.1)]{drenth2025thesis},
{\small
\setlength{\arraycolsep}{2.5pt}
\begin{equation}
  \begin{bmatrix}\dot q\\ \ddot q\\ \hline z\\ \hline \qact\end{bmatrix}
  =\underbrace{\left[\begin{array}{cc|cc|c}
  0&I_3&0&0&0\\
  -RK&-RC&-RM_1&-RM_2&RP\\ \hline
  -RK&-RC&-RM_1&-RM_2&RP\\
  0&0&I_3&0&0\\ \hline
  P^\top&0&0&0&0
  \end{array}\right]}_{G}
  \begin{bmatrix}q\\ \dot q\\ \hline w\\ \hline \uact\end{bmatrix},
  \label{eq:lfr_G}
\end{equation}
}
where $R=M_0^{-1}$, $\tilde x=\col(q,\dot q)\in\R^6$ is the state, and
$G\in\R^{15\times15}$.
$G$ depends on the physical parameters but not on $Y$.

Closing the loop with \eqref{eq:delta} gives $a_1=Y\ddot q$ and
$a_2=Y^2\ddot q$.
Substituted in \eqref{eq:lfr_constant}, these return \eqref{eq:eom}, so the
LFR realises the baseline exactly.
The LFR is well posed if $I_6-D_{zw}\Dsch(Y)$ is nonsingular for every
$Y$~\cite[Sec.~2.1]{drenth2025thesis}, where $D_{zw}$ is the block of $G$
from $w$ to $z$.
For our realisation, the Schur complement of its lower identity block gives
(Appendix~\ref{app:lfr-details})
\begin{equation}
  \det\bigl(I_6-D_{zw}\Dsch(Y)\bigr)=\frac{\det M(Y)}{\det M_0},
  \label{eq:wellposed}
\end{equation}
so the LFR is well posed exactly where $M(Y)$ is nonsingular.

Positive masses and inertias give $M(Y)\succ0$ for every $Y\in\R$
(Appendix~\ref{app:lfr-wellposed}), so the baseline is well posed at every
payload position, not only on the stroke.
The implementation therefore solves the loop in closed form, with
$M(Y)^{-1}=\operatorname{adj}M(Y)/\det M(Y)$
(Appendix~\ref{app:lfr-details}).
It evaluates $G$ with the resulting $w$.
The LFR remains the form in which LPV-LFR model augmentation takes its
baseline~\cite[Sec.~5.2, Eq.~(5.1)]{drenth2025thesis}.
It also yields the exact well-posedness condition \eqref{eq:wellposed}.

Two realisation choices fix the time domain and the size of the Delta-block.
Unlike the discrete-time LPV-LFR of Drenth et
al.~\cite[Sec.~3.1]{drenth2025lpvlfr}, it stays in continuous time, where
$\theta$ enters through $M$, $C$ and $K$ as derived.
Section~\ref{sec:aug_structure} discretises it with one RK4 step per sample.
We do not claim this six-channel realisation minimal.
\todo{Missing: reason for the six-channel realisation rather than a reduced one (header ``Choices to justify''; the six-to-four reduction is excluded from the text by the meeting audit). Candidate: every block of $G$ in \eqref{eq:lfr_G} is $M_0^{-1}$ times a physical matrix, so $\theta$ stays readable in $G$; a reduced realisation mixes the parameters (own reading, no decision entry).}

\subsection{Physical parameters and identifiable combinations}\label{sec:params}
\ifdraft
\begin{itemize}
\item The input-output behaviour fixes $M_0,M_1,M_2,C,K$ and conversely, so the identifiable part of $\theta$ is what the matrices contain (this work)
\item Twelve nonzero entries, three equal to $\pm m_h$, leave ten independent functions: exactly ten combinations are identifiable (this work; definition Ovchinnikov et al.\ Def.~7)
\item Written physically as $\theta_{\mathrm{base}}$; four lost directions; estimates reported in $\theta_{\mathrm{base}}$ (this work; Gautier and Khalil counterpart unverified (?))
\item For $m_h,m_\Sigma,J_{\mathrm{eff}}>0$ the combinations do not give $M(Y)\succ0$; that holds iff \eqref{eq:lfr_admissibility} (this work; proof Appendix~\ref{app:lfr-wellposed})
\item Training keeps \eqref{eq:lfr_admissibility} at every iterate; the map and its justification are in Section~\ref{sec:aug_closed_loop}
\end{itemize}
\fi
The fourteen parameters enter the model only through $M_0$, $M_1$, $M_2$,
$C$ and $K$.
Two parameter vectors with the same input-output behaviour give the same $q$
for every input, since $P$ is invertible.
The difference of their equations of motion then vanishes on every
trajectory.
The gantry is fully actuated, so every triple $(q,\dot q,\ddot q)$ occurs on
some trajectory.
The two vectors therefore give the same five matrices.
Equal matrices give the same model, so the identifiable part of $\theta$ is
what these matrices contain.

Up to symmetry, \eqref{eq:mass_matrix} and \eqref{eq:damping_stiffness_matrices}
have twelve nonzero entries.
Three of them, $M_{0,33}$, $-M_{1,12}$ and $M_{2,22}$, equal $m_h$, which
leaves ten distinct functions of $\theta$.
Their Jacobian with respect to $\theta$ has rank ten.
Exactly ten independent functions of $\theta$ are therefore
identifiable~\cite[Def.~7]{ovchinnikov2021computing}.
We report estimates in these ten combinations, written in physically
readable form as
\begin{equation}
  \theta_{\mathrm{base}}=[k_{b,\Sigma},c_{g1},c_{g2},c_y,c_{b,\Sigma},
  m_h,m_\Sigma,m_\Delta,J_{\mathrm{eff}},d]^\top\in\R^{10},
  \label{eq:phi}
\end{equation}
where $k_{b,\Sigma}=k_{b1}+k_{b2}$, $c_{b,\Sigma}=c_{b1}+c_{b2}$,
$m_\Sigma=m_1+m_2+m_b$, $m_\Delta=m_1-m_2$ and
$J_{\mathrm{eff}}=J_b+J_h+\tfrac{L_b^2}{4}(m_1+m_2)$.
The four lost directions are $k_{b1}-k_{b2}$, $c_{b1}-c_{b2}$, $J_b-J_h$, and
a transfer of mass from $m_1$ and $m_2$ to $m_b$ with a compensating change
of $J_b+J_h$.
\todo{Missing: whether to name $\theta_{\mathrm{base}}$ the gantry counterpart of the base inertial parameters of robotics (Gautier and Khalil 1990, key \texttt{gautier1990minimum}). Candidate: add the clause and the citation after reading the paper; no local PDF, unverified (README source map row II-C).}
\todo{Missing: notation. The index $\Delta$ in $m_\Delta$ clashes with the scheduling block $\Dsch$, and $m_\Delta$ is used in Sections~\ref{sec:aug_closed_loop} and~\ref{sec:obc}. Candidate: $m_{\mathrm{diff}}$, the code name (Dirk decides, all three sections change together).}

In $\theta_{\mathrm{base}}$, positive definiteness of the inertia is no longer
automatic, because $m_\Delta$ can take values that no positive masses
produce.
For $m_h,m_\Sigma,J_{\mathrm{eff}}>0$, $M(Y)\succ0$ for every $Y\in\R$ if
and only if
\begin{equation}
  m_\Sigma J_{\mathrm{eff}}>\tfrac{L_b^2}{4}\,m_\Delta^2
  \label{eq:lfr_admissibility}
\end{equation}
(Appendix~\ref{app:lfr-wellposed}).
Section~\ref{sec:aug_closed_loop} therefore trains $\theta_{\mathrm{base}}$
in coordinates that keep \eqref{eq:lfr_admissibility}, and with it
well-posedness, at every iterate.

\subsection{Compared baselines}\label{sec:baselines}
\ifdraft
\begin{itemize}
\item $\mathcal M_{\mathrm{LPV}}$ and $\mathcal M_{\mathrm{LTI}}$ map $\uact$ to $\qact$; neither is trained; both use the parameters of the data-generating system (THESIS-RESULTS step~1)
\item $\mathcal M_{\mathrm{LTI}}$ evaluates the Delta-block at a fixed $Y_{\mathrm{op}}$, so it is LTI; it measures what the scheduling contributes (Aspect~1; D-242); value of $Y_{\mathrm{op}}$ in Section~V
\end{itemize}
\fi
Section~\ref{sec:results} compares two baselines to measure what the
scheduling contributes (Aspect~1 of Section~I).
Neither is trained.
With the parameters of the data-generating system and the payload position
$Y_{\mathrm{op}}$ of Section~\ref{sec:setup}, they are
\begin{equation}
\begin{aligned}
  \mathcal M_{\mathrm{LPV}}&:\quad\eqref{eq:lfr_G},\quad w=\Dsch(Y)z,\\
  \mathcal M_{\mathrm{LTI}}&:\quad\eqref{eq:lfr_G},\quad w=\Dsch(Y_{\mathrm{op}})z,
\end{aligned}
\label{eq:baselines}
\end{equation}
where $Y_{\mathrm{op}}\in\R$ is fixed.
Closing the loop of $\mathcal M_{\mathrm{LTI}}$ gives \eqref{eq:eom} with the
constant inertia $M(Y_{\mathrm{op}})$, an LTI model.
\todo{Missing: model notation. Candidate: $\mathcal M_{\mathrm{LPV}}$ and $\mathcal M_{\mathrm{LTI}}$, since $M$ is the inertia matrix (README open conflict on model names; Section~\ref{sec:aug_closed_loop} uses $\mathcal M_{\mathrm U}$).}
\todo{Missing: no runner evaluates $\mathcal M_{\mathrm{LTI}}$ yet; the frozen-$Y$ option of the baseline block is a capability only (README source map, II Frozen-LTI baseline). Candidate: add the step~1 evaluation of THESIS-RESULTS before Section~\ref{sec:results} reports it.}
